\documentclass[output=paper]{langscibook}
\ChapterDOI{10.5281/zenodo.22880294}
\title{\textit{Vavaaiku} – Angling with a swinging bait}
\author{Enoch Horai Magum and Shalom Magum}
\abstract{}
\begin{document}
\maketitle\label{chap:33}
\begin{xltabular}{\textwidth}{lXX}
	(\ref{ex:text26-1})&\textit{O kaku vagana to dao raara bona vavaaiku na pakupaku raara namana, beara hae batara teo sinivi.}&The kind of fishing that we call \textit{vavaaiku} `letting it swing', we do it in the ocean while we are sitting in the canoe.\\
	(\ref{ex:text26-2})&\textit{Na pakupaku raara bona kanono ae a iri ae sa taba sa ponisi mene dee gunaha bona kanono.}&We do it with a string and a hook and something heavy to get the string down.\\
	(\ref{ex:text26-3})&\textit{Beara mei rara tea tahoa, amaa taba vaarii toro vaamatatopo.}&Before we set off, these things must be prepared.\\
	(\ref{ex:text26-4})&\textit{A kanono to vavaaiku ni raara na paku vonaen:}&The string that we use for angling is made as follows:\\
	(\ref{ex:text26-5})&\textit{Eara na kisi ni roho rara a ta booru tea kanono, eara repaa kisi mana nia iri putaa.}&We first tie a piece of lead to the string and then we tie the hook below it.\\
	(\ref{ex:text26-6})&\textit{Ge beara rake, eara repaa kisi nia iri gina bea ta booru en putaa.}&Or if we want, we tie the hook on top while the piece of lead is below the hook.\\
	(\ref{ex:text26-7})&\textit{Teo kaku vagana vai, eara na beana rarai.}&For this kind of fishing, we bait it (the hook).\\
	(\ref{ex:text26-8})&\textit{Eara he toro beana nio urita, iobo ge a iana komana vahaa, beara peepee varutarutaa e.}&We must bait it with an octopus, a sea-saussage, or even a fish when we have cut it up into little pieces.\\
	(\ref{ex:text26-9})&\textit{Bea maamihu taba vai matatopo, eara repaa hua hiava namana.}&When everything is prepared, we paddle up to the deep sea.\\
	(\ref{ex:text26-10})&\textit{Eara repaa vaagunaha a kanono, vaantee koa tobina, eara repaa vaabunobuno koa ni rae, a beana re potee voen na vavaaiku nana.}&Then we let the string down, but only half way, and then move it to and fro so that the bait looks like swinging.\\
	(\ref{ex:text26-11})&\textit{Teebona to kahi avuhu ae tara mau o iana bona, eori repaa ma agomo bona beana.}&So the fish will smell and see it, and then they come and swallow the bait.\\
	(\ref{ex:text26-12})&\textit{Beara haraa voen na rahi batana bona kanono, eara repaa vaatasu hiava nia kanono,}&When we feel that a fish is pulling the string, we fling the string upwards,\\
	(\ref{ex:text26-13})&\textit{a iri repaa kao hopo rivona iana, eara repaa rahi hiava e kou.}&so that the hook moves into the mouth of the fish, and then we pull it up.\\
	(\ref{ex:text26-14})&\textit{Eara na antee rara tea paku o kaku vagana vai tea bon, tea roava, teo tarahi ge teo puhana.}&We can do this kind of fishing at night, during the day, when the tide is ebbing or at high tide.\\
\end{xltabular}
\end{document}